\documentclass[aps,prl,onecolumn,amsmath,amssymb]{revtex4-2}

\usepackage{bm}
\usepackage{booktabs}
\usepackage{xcolor}
\usepackage[colorlinks=true,citecolor=blue!55!black,urlcolor=blue!55!black]{hyperref}

\newcommand{\Jzz}{J_{zz}}
\newcommand{\Jpm}{J_{\pm}}
\newcommand{\bq}{\bm q}
\newcommand{\btheta}{\bm\theta}
\newcommand{\PP}{\mathcal P}

\begin{document}

\title{Supplemental Material for ``Transport-projected exact diagonalization
of quantum spin ice: benchmarks and convergence criteria''}
\author{Z.-B. Zhou}
\affiliation{(affiliation)}
\date{\today}
\maketitle

\section{Model, convention, and cluster geometry}

We use local spin axes and
\begin{equation}
H=H_0+V=\Jzz\sum_{\langle ij\rangle}S_i^zS_j^z
-\Jpm\sum_{\langle ij\rangle}(S_i^+S_j^-+S_i^-S_j^+).
\label{eq:model_si}
\end{equation}
Huang \emph{et al.} use this same convention at the sign-free point
$\Jpm/\Jzz=0.046$~\cite{Huang2018}.  No sign or factor-of-two conversion is
therefore needed for the external benchmark.

Sites are represented in a periodic covering lattice.  For each nearest
neighbor pair $(i,j)$, the stored integer image $\bm n_{ij}$ gives the lifted
bond displacement
\begin{equation}
 \bm d_{ij}=\bm r_j-\bm r_i-\bm n_{ij}^{T}L,
\end{equation}
where the rows of $L$ are the three supercell vectors.  A depth-first cycle
enumeration retains simple cycles that preserve the ice constraint and, for
six-cycles, are planar in the covering lattice.  The resulting census is
\begin{center}
\begin{tabular}{lrrrr}
\toprule
cluster & ice states & winding 4 & contractible 6 & wrapping 6\\
\midrule
cubic-16 & 90 & 36 & 16 & 48\\
FCC-32 & 2970 & 24 & 32 & 96\\
\bottomrule
\end{tabular}
\end{center}
There are no contractible four-cycles.  The active geometry module performs
the enumeration independently from the archived campaigns.

\section{Canonical block transformation}

Let $P$ project onto the degenerate ice manifold of $H_0$ and $Q=1-P$.
Introduce a bookkeeping parameter $\lambda$ multiplying $V$.  We seek
\begin{equation}
 S^{(N)}=\sum_{n=1}^{N}\lambda^n S_n,
 \qquad S_n^\dagger=-S_n,
 \qquad PS_nP=QS_nQ=0,
\label{eq:generator}
\end{equation}
such that the off-diagonal block of
$\widetilde H^{(N)}=e^{-S^{(N)}}He^{S^{(N)}}$ vanishes through order $N$.
The last two conditions fix the canonical gauge~\cite{Bravyi2011}.

For completeness, define $\mathcal L_0(X)=[H_0,X]$.  Suppose orders below
$n$ have been fixed and let $B_n$ be the coefficient of $\lambda^n$ in the
off-diagonal part of the Baker--Campbell--Hausdorff series with $S_n$ omitted.
Then
\begin{equation}
 S_n=-\mathcal L_0^{-1}(B_n),
\label{eq:recursion}
\end{equation}
where the inverse acts only between $P$ and $Q$.  It exists while the Ising
charge gap separating the chosen blocks remains nonzero.  Equations
(\ref{eq:generator}) and (\ref{eq:recursion}) recursively define the series
to arbitrary order; they are not a third-order ansatz.

The two finite-order blocks are
\begin{align}
 H_P^{(N)}&=P\widetilde H^{(N)}P,\\
 H_Q^{(N)}&=Q\widetilde H^{(N)}Q,
\end{align}
and the omitted off-diagonal block is $O(\lambda^{N+1})$.  An $N$-order
calculation is therefore used only after comparing with $N+1$ (or an exact
band) over the temperature and frequency window reported.

At second and third order the row form used for the topology diagnostic is
\begin{align}
 H_P^{(2)}&=PVRVP,\\
 H_P^{(3)}&=PVRVRVP,
 \qquad R=-Q(H_0-E_{\rm ice})^{-1}Q.
\end{align}
The code stores each virtual sequence before rows with equal source, target,
and image transport are merged.  This is what permits a transport projection
without inferring path content from the final matrix element.

\section{Transported dipole and character projection}

Consider a completed sequence that maps ice state $|a\rangle$ to
$|b\rangle$.  Every elementary $S_i^+S_j^-$ transports one spin unit from
$j$ to $i$.  Its lifted displacement is accumulated along the virtual
sequence.  In integer supercell coordinates the total is $\bq\in\mathbb Z^3$.
Equivalently, in Cartesian coordinates one may combine the home-cell dipole
change and the accumulated bond images,
\begin{equation}
 \bm\delta_{ba}=\sum_{i\in\mathrm{raised}}\bm r_i
 -\sum_{j\in\mathrm{lowered}}\bm r_j+\bm N^TL.
\label{eq:delta}
\end{equation}
Changing a site's home-cell representative shifts the two terms oppositely,
so $\bm\delta$ and the statement $\bq=0$ are gauge invariant.

Attach a phase $e^{i\btheta\cdot\Delta\bq}$ to every elementary hop, where
$\Delta\bq$ is its lifted source increment.  Individual increments may be
half-integer, while every completed ice-to-ice path has integer $\bq$.  A
completed row then
has the Fourier expansion
\begin{equation}
 [H_P^{(N)}(\btheta)]_{ba}
 =\sum_{\bq}h_{ba}^{(N)}(\bq)e^{i\btheta\cdot\bq}.
\label{eq:fourier}
\end{equation}
In the active exact calculation the primitive completed-path source is
$\bq=2\bm\delta$; the quarter-cell integer used by the row code is therefore
$\bm\delta_4=2\bq$.  This single normalization is used at every $M$.
The discrete orthogonality identity
\begin{equation}
 \frac{1}{M}\sum_{m=0}^{M-1}e^{2\pi i m q/M}
 =\begin{cases}1,&q=0\pmod M,\\0,&\text{otherwise},\end{cases}
\end{equation}
proves
\begin{equation}
 \PP_M[H_P^{(N)}]
 =\sum_{\bq=0\ ({\rm mod}\ M)}h^{(N)}(\bq).
\label{eq:alias}
\end{equation}
For fixed $N$ only finitely many hops occur.  Choosing $M$ larger than every
component of their possible net transport makes Eq.~(\ref{eq:alias}) the
exact $\bq=0$ selector.  At $N\to\infty$ arbitrarily large winding is
possible, so a finite grid aliases multiple winding and $M$ convergence is
mandatory.

This explains two rejected alternatives.  Averaging
$\mathrm{Tr}\,e^{-\beta H(\btheta)}$ is nonlinear and does not define the
partition function of the projected operator.  Averaging microscopic
Hamiltonians projects individual hops before they compose and therefore does
not classify closed virtual paths.  A fitted counterterm can cancel one
leading matrix element, but it does not implement Eq.~(\ref{eq:alias}) at
higher order or all temperatures.

\section{Exact isolated-band limit}

Let $P_\lambda$ project onto an isolated exact band with
$\mathrm{rank}(P_\lambda)=\mathrm{rank}(P)$.  If $PP_\lambda P$ is positive
definite on $P$, the partial isometry
\begin{equation}
 W_\lambda=P_\lambda P(PP_\lambda P)^{-1/2}
\label{eq:polar}
\end{equation}
satisfies $W_\lambda^\dagger W_\lambda=P$ and
$W_\lambda W_\lambda^\dagger=P_\lambda$.  The pulled-back operator
\begin{equation}
 H_{\rm band}=W_\lambda^\dagger H W_\lambda
\label{eq:band}
\end{equation}
is the canonical all-order continuation of $H_P^{(N)}$ while the band remains
isolated.  Numerically, if $\Psi$ contains the exact band eigenvectors and
$X=P\Psi$, Eq.~(\ref{eq:polar}) becomes the Loewdin orthogonalization
$X(X^\dagger X)^{-1/2}$.  The minimum eigenvalue of $X^\dagger X$ is a
required overlap diagnostic.  A closing separation from the complement or a
singular Gram matrix invalidates the selected-band protocol.

For every source angle $\btheta$, the same polar convention must be used
before character averaging.  Changing the arbitrary eigenvector gauge at
each angle would otherwise produce a meaningless operator average.  The
active exact-band comparison uses this convention.  At
$\Jpm/\Jzz=0.046$, the centered changes are 6.78\% from $M=2$ to 3 and
0.498\% from $M=3$ to 4.  The latter passes the 5\% character-convergence
gate.  The zero-source minimum ice overlap is 0.762, and every generic
$M=3,4$ point has a larger overlap.

\section{Full-temperature thermodynamics}

Let $\{E_i\}$ be the full microscopic spectrum, $B$ the indices of the bare
selected band, and $\{\widetilde E_a\}$ the clean-band spectrum.  Projection
preserves the band trace; for a truncated series we enforce the same trace to
fix the band's energy relative to the complement.  Define, for $k=0,1,2$,
\begin{equation}
 M_k^{\rm clean}(\beta)=
 \sum_{i\notin B}E_i^k e^{-\beta E_i}
 +\sum_a\widetilde E_a^k e^{-\beta\widetilde E_a}.
\label{eq:moments}
\end{equation}
Then
\begin{align}
 E&=M_1/M_0,\\
 C&=\beta^2\left(M_2/M_0-(M_1/M_0)^2\right),\\
 S&=\ln M_0+\beta E.
\end{align}
Equation~(\ref{eq:moments}) is equivalent to diagonalizing one Hermitian
operator whose selected block is replaced and whose complement is unchanged.
For cubic-16 we use the complete 65,536-level spectrum.  For FCC-32 the same
identity requires stochastic trace estimation of the complement; that
calculation has not yet been performed.

\section{QMC benchmark provenance and metrics}

The reference heat capacity and entropy are vector extractions from
Fig.~1(b) of Huang \emph{et al.}~\cite{Huang2018}.  The arXiv source contains
the original vector PDF, allowing the red $C/N$ and orange $S/N$ paths to be
read without raster tracing.  The extraction script, axis calibration, CSV,
and provenance are stored under \texttt{campaign/data/}.  Estimated
extraction uncertainties are 0.002 in $\log_{10}(T/\Jzz)$, 0.0015 in $C/N$,
and 0.002 in $S/N$.  These are not substitutes for unreported QMC statistical
errors.

Curve error is the root-mean-square difference after interpolation in
$\log T$ over the digitized support.  Thresholds are defined in
\texttt{campaign/config.json} before evaluation.  The current values are
\begin{center}
\begin{tabular}{lccc}
\toprule
quantity & bare cubic-16 & clean cubic-16 & QMC\\
\midrule
low $T_{\rm p}/\Jzz$ & 0.02554 & 0.001463 & 0.001100\\
low-$T$ $C/N$ RMSE & 0.0859 & 0.02327 & 0\\
high $T_{\rm p}/\Jzz$ & 0.2549 & 0.2549 & 0.2105\\
$S/N$ RMSE & 0.0514 & 0.04509 & 0\\
\bottomrule
\end{tabular}
\end{center}
The converged $M=4$ result reduces the low-temperature heat-capacity RMSE by
72.9\%, but its final value $0.02327$ remains above the preregistered 0.02
tolerance.  The entropy gate also fails because $0.04509>0.02$.  The
discrepancy is concentrated in the low-temperature region, where the finite
clean band retains excess entropy relative to QMC.

The same paper specifies a future dynamics gate.  It uses an $8\times8\times8$
primitive-cell lattice and reports
\begin{equation}
 S^{zz}(\bm q,\omega)=\sum_{\alpha=1}^4S^{zz}_{\alpha\alpha}(\bm q,\omega),
 \qquad
 S^{+-}(\bm q,\omega)=\sum_{\alpha=1}^4S^{+-}_{\alpha\alpha}(\bm q,\omega),
\end{equation}
at $T/\Jzz=0.001$, 0.04, and 0.1 along
$2\pi(0,2k,0)/8$ and $2\pi(k,k,k)/8$~\cite{Huang2018}.  ED comparison must
use only momenta realized by its finite translation group, the same
sublattice trace and spin normalization, a linear frequency axis, and a
stated broadening comparable to the QMC-SAC resolution.  This matched
dynamics calculation has not yet been performed.

\section{Complete validation status}

\begin{center}
\begin{tabular}{lll}
\toprule
gate & status & evidence or missing calculation\\
\midrule
loop topology & pass & winding 4/6 removed; contractible 6 retained\\
all-temperature stability & pass & cubic-16 high-$T$ peak position unchanged\\
QMC heat capacity & fail & exact $M=4$ low-$T$ RMSE $0.02327>0.02$\\
QMC entropy & fail & exact $M=4$ RMSE $0.04509>0.02$\\
character resolution & pass & centered $M=3\to4$ error 0.498\%\\
band overlap & pass & source-resolved minimum is 0.762\\
order convergence & diagnostic only & production spectrum is exact-band\\
QMC-SAC dynamics & pending & matched $S^{zz}$ and $S^{+-}$ required\\
QMC energy & pending & not plotted in the benchmark paper\\
FCC-32 full Hamiltonian & pending & current result is order-three gauge block\\
\bottomrule
\end{tabular}
\end{center}

The next production sequence is therefore fixed: include all conserved
magnetization sectors in the unrestricted gap diagnostic; resolve the failed
heat and entropy tests through cluster scaling; match the Huang \emph{et al.}
QMC-SAC spectra; and only then run the full-Hamiltonian FCC-32 and $\pi$-flux
campaigns.

\bibliography{references}

\end{document}
